%%%%%
%
%% Dans ce fichier, les �ventuels accents sont cod�s
%% en ISO-8859-15 (aka latin9)
%% La ligne suivante donne trois exemples 
%% e accent aigu : �, e accent grave : �, e-dans-l'o : 1/2
\section[Rapport de m�lange en vapeur d'eau � partir du
point de ros�e ou de givre]{Rapport de m�lange en vapeur d'eau � partir de la temp�rature du
point de ros�e ou de givre}


\itemstyle{Op�rateur~:} \operfmt{RapM\_Tdf}
\label{OP-RapM_Tdf}
\index{RapM\_Tdf}

%%%%%OBJET%%%%
\itemstyle{Objet~:}
Calcul du rapport de m�lange en vapeur d'eau � partir
de la temp�rature du point de ros�e ou de givre et 
de la pression de l'air.

%%%%%ARGUMENTS%%%%
\bigskip
%% Dans ce fichier, les �ventuels accents sont cod�s
%% en ISO-8859-15 (aka latin9)
%% La ligne suivante donne trois exemples 
%% e accent aigu : �, e accent grave : �, e-dans-l'o : �
%%%%%ENTREES%%%%
\begin{tabarg}
\multicolumn{3}{l}{\titretab{Entr\'ees}}\\
\hline
$T_{df}$ & Vecteur, \texttt{float} & Temp�rature du point de ros�e ou de givre ($K$ ou $^\circ C$)\\
$P_s$ & Vecteur, \texttt{float} & Pression statique ($hPa$) \\
\end{tabarg}

\medskip
%%%%%SORTIES%%%%
\begin{tabarg}
\multicolumn{3}{l}{\titretab{Sorties}}\\
\hline
$RapM$ & Vecteur, \texttt{float} & Rapport de m�lange en vapeur d'eau (g/kg)\\
\end{tabarg}


%%%%%FORMULE%%%%
\itemstyle{Formule}

%% Formule de l'operateur : RapM_Tdf
%% Dans ce fichier, les �ventuels accents sont cod�s
%% en ISO-8859-15 (aka latin9)
%% La ligne suivante donne trois exemples 
%% e accent aigu : �, e accent grave : �, e-dans-l'o : �
\[
RapM = \frac{R_a}{R_v} . \frac{e}{P_s - e } \qquad \text{ avec }
e = 
\begin{cases}
 e_w(T_{df}) \text{ si } T_{df} \ge 0.0 \text{�C}\\
 e_{wi}(T_{df}) \text{ sinon}\\
\end{cases}
\]
et $\frac{R_a}{R_v} = \frac{M_v}{M_a} = 0.622$


%%%%COMMENTAIRE%%%%
\itemstyle{Commentaire~:}

Le fonctions $e_{w}(t)$ et $e_{wi}(t)$ sont d�crites en annexe
\ref{annexe-thermo}.


%%%%%AUTEUR%%%%
\itemstyle{Auteur~:}
CNRM/GMEI/TRAMM

%%%%%REFERENCE%%%%
\itemstyle{R�f�rences~:}

Triplet-Roche\cite{Triplet-Roche}, p55\\
et d�finitions des humidit�s (relatives, absolues, sp�cifiques)
et du rapport de m�lange
selon la table 4.3 des tables m�t�orologiques internationales
de l'OMM \cite{OMM-188.TP.94}.
